% concatenated from vadic2.tex, chapter1/1motivation.tex, chapter1/2summary.tex, chapter1/3notation.tex, chapter2/backcm1.tex, chapter2/backcm2.tex, chapter2/backcm3.tex, chapter2/backcm4.tex, chapter3/background1.tex, chapter3/grosszagier.tex, chapter4/crysbasis1.tex, chapter4/crysbasis2.tex, chapter4/crysbasis3.tex, chapter4/crysbasis4.tex, chapter5/cmapplication1.tex, chapter5/cmapplication2.tex, chapter5/cmapplication4.tex
\documentclass[12pt,twoside,a4paper]{article}
\usepackage[english]{babel}%Packages
\usepackage[utf8]{inputenc}
\usepackage{amsmath}
\usepackage{amsfonts}
\usepackage[pagebackref]{hyperref}
\usepackage{amsthm}
\usepackage{thmtools}
\usepackage{amssymb, graphics, setspace}
\usepackage{bm}
\usepackage{tikz-cd}
\usepackage{xr}
\usepackage[nameinlink]{cleveref}
\usepackage{url}
\usepackage{xcolor}

\usepackage{enumitem}
\usepackage{authblk}
\usepackage[title]{appendix}
\usepackage{chngcntr}
\usepackage{apptools}
\usepackage{mathtools}
\usepackage{bookmark}
\usepackage{appendix}
\definecolor{crimsonglory}{rgb}{0.75, 0.0, 0.2}
\definecolor{darkpowderblue}{rgb}{0.0, 0.2, 0.6}
\hypersetup{
	colorlinks   = true, %Colours links instead of ugly boxes
	urlcolor     = darkpowderblue, %Colour for external hyperlinks
	linkcolor    = darkpowderblue, %Colour of internal links
	citecolor   = crimsonglory %Colour of citations
}


\newcommand{\subscript}[2]{$#1 _ #2$}%Changing enumeration styles.
\theoremstyle{plain}%Theorem numbering stuff.
\newtheorem{theorem}{Theorem}[section] % reset theorem numbering for each chapter

\newtheorem{definition}[theorem]{Definition} % definition numbers are dependent on theorem numbers
\newtheorem{exmp}[theorem]{Example} % same for example numbers
\newtheorem{conv}[theorem]{Convention} % definition numbers are dependent on theorem numbers
\newtheorem{prop}[theorem]{Proposition}
\newtheorem{cor}[theorem]{Corollary}
\newtheorem{lemma}[theorem]{Lemma}
\newtheorem{claim}[theorem]{Claim}
\newtheorem{remark}[theorem]{Remark}
\newtheorem{remarks}[theorem]{Remarks}
\newtheorem{conj}[theorem]{Conjecture}
\newtheorem{assumption}[theorem]{Assumption}
\newtheorem{para}[theorem]{Paradigm}

\DeclareMathOperator{\card}{card}%Stuff that are supposed to make my life easier.
\DeclareMathOperator{\res}{res}
\DeclareMathOperator{\id}{id}
\DeclareMathOperator{\homm}{Hom}
\DeclareMathOperator{\Res}{Res}
\DeclareMathOperator{\Inf}{Inf}
\DeclareMathOperator{\Cor}{Cor}
\DeclareMathOperator{\ver}{Ver}
\DeclareMathOperator{\End}{End}
\DeclareMathOperator{\GL}{GL}
\DeclareMathOperator{\SL}{SL}
\DeclareMathOperator{\coker}{Coker}
\DeclareMathOperator{\tors}{tors}
\DeclareMathOperator{\diag}{diag}
\DeclareMathOperator{\rank}{rank}
\DeclareMathOperator{\pr}{pr}
\DeclareMathOperator{\tr}{tr}
\DeclareMathOperator{\cl}{cl}
\DeclareMathOperator{\spec}{Spec}
\DeclareMathOperator{\gal}{Gal}
\DeclareMathOperator{\aut}{Aut}
\DeclareMathOperator{\gr}{Gr}
\DeclareMathOperator{\inv}{inv}
\DeclareMathOperator{\ram}{Ram}
\DeclareMathOperator{\Span}{Span}
\DeclareMathOperator{\LT}{LT}
\DeclareMathOperator{\fin}{fin}
\DeclareMathOperator{\atyp}{atyp}
\DeclareMathOperator{\disc}{disc}
\DeclareMathOperator{\zar}{Zar}
\DeclareMathOperator{\chara}{char}
\DeclareMathOperator{\crys}{crys}
\DeclareMathOperator{\isog}{isog}
\DeclareMathOperator{\ord}{ord}
\DeclareMathOperator{\SP}{SP}
\DeclareMathOperator{\adj}{adj}
\DeclareMathOperator{\arch}{arch}
\DeclareMathOperator{\new}{new}
\DeclareMathOperator{\ssing}{ssing}

\newcommand{\im}[1]{Im(#1)}
\newcommand{\Sum}[2]{\displaystyle\sum_{#1}^{#2}}
\newcommand{\Bigcup}[2]{\displaystyle\cup_{#1}^{#2}}
\newcommand{\Bigsum}[2]{\displaystyle\bigoplus_{#1}^{#2}}
\newcommand{\tensor}[2]{\displaystyle\otimes_{#1}^{#2}}
\newcommand{\Z}{\mathbb{Z}}
\newcommand{\N}{\mathbb{N}}
\newcommand{\C}{\mathbb{C}}
\newcommand{\Q}{\mathbb{Q}}
\newcommand{\B}{\mathbb{B}}
\newcommand{\A}{\mathbb{A}}
\newcommand{\R}{\mathbb{R}}
\newcommand{\F}{\mathbb{F}}
\newcommand{\G}{\mathbb{G}}
\newcommand{\fgj}{\hat{\gamma}^{*}_j}
\newcommand{\drfj}[1]{\hat{\gamma}^{*}_{{#1},B_{dR}}}
\newcommand{\CE}{\mathcal{E}}
\newcommand{\CX}{\mathcal{X}}
\newcommand{\ld}{,\ldots,}
\newcommand{\CP}{\mathcal{P}}
\newcommand{\CO}{\mathcal{O}}
\newcommand{\unr}{unr}
\newcommand{\CF}{\mathcal{F}}
\newcommand{\et}{\acute{e}t}
\newcommand{\proet}{pro\acute{e}t}
\newcommand{\HL}{HL}
\newcommand{\V}{\mathbb{V}}
\usepackage[OT2,T1]{fontenc}\DeclareFontSubstitution{OT2}{wncyr}{m}{n}% original Sha encoding
\DeclareSymbolFont{cyrletters}{OT2}{wncyr}{m}{n}
\DeclareMathSymbol{\Sha}{\mathalpha}{cyrletters}{"58}
\usepackage{tocloft}%Table of Contents
\renewcommand\cftsecfont{\small}
\renewcommand\cftsecpagefont{\small}
\renewcommand\cftsubsecfont{\small}
\renewcommand\cftsubsecpagefont{\small}

%For matrices with brakcets.
\newcommand\coolover[2]{\mathrlap{\smash{\overbrace{\phantom{%
					\begin{matrix} #2 \end{matrix}}}^{\mbox{$#1$}}}}#2}

\newcommand\coolunder[2]{\mathrlap{\smash{\underbrace{\phantom{%
					\begin{matrix} #2 \end{matrix}}}_{\mbox{$#1$}}}}#2}

\newcommand\coolleftbrace[2]{%
	#1\left\{\vphantom{\begin{matrix} #2 \end{matrix}}\right.}

\newcommand\coolrightbrace[2]{%
	\left.\vphantom{\begin{matrix} #1 \end{matrix}}\right\}#2}

\makeatletter
\newcommand{\address}[1]{\gdef\@address{#1}}
\newcommand{\email}[1]{\gdef\@email{\url{#1}}}
\newcommand{\@endstuff}{\par\vspace{\baselineskip}\noindent\small
	\begin{tabular}{@{}l}\scshape\@address\\\textit{E-mail address:} \@email\end{tabular}}
\AtEndDocument{\@endstuff}
\makeatother

\title{On the $v$-adic values of G-functions II:\\
	\Large Towards Effective Brauer-Siegel}
	\author{Georgios Papas}
\address{Faculty of Mathematics and Computer Science\\
	The Weizmann Institute of Science\\
	234 Herzl Street, Rehovot 76100, Israel, and\\\\	
	
	Institute for Advanced Study\\
	1 Einstein Drive\\
	Princeton, N.J. 08540\\
	U.S.A.}
\email{georgios.papas@weizmann.ac.il, gpapas@ias.edu}

\newcommand{\mathsym}[1]{{}}
\newcommand{\unicode}[1]{{}}
